\begin{figure}[t]
\centering
\resizebox{\linewidth}{!}{%
\begin{tikzpicture}[
  font=\small,
  node/.style={circle, draw=black!48, line width=0.55pt, align=center, minimum size=0.72cm, inner sep=1pt, fill=white},
  qnode/.style={node, fill=blue!10, minimum size=0.86cm},
  pnode/.style={node, fill=cyan!9},
  dnode/.style={node, fill=gray!10},
  enode/.style={node, fill=orange!13},
  mnode/.style={node, fill=green!12},
  rnode/.style={node, fill=purple!9},
  hedge/.style={draw, rounded corners=8pt, line width=0.7pt, fill opacity=0.18},
  edge/.style={draw=black!45, line width=0.55pt},
  arrow/.style={-{Latex[length=1.8mm]}, draw=black!55, line width=0.55pt},
  legend/.style={draw=black!20, rounded corners=2pt, fill=black!2, align=left, inner sep=4pt}
]
\node[qnode] (q) at (0,0.5) {$q$};
\node[pnode] (p1) at (2.3,1.45) {$p_1$};
\node[pnode] (p2) at (2.25,-0.35) {$p_2$};
\node[pnode] (p3) at (4.1,0.55) {$p_3$};
\node[dnode] (d1) at (5.9,1.45) {$d_A$};
\node[dnode] (d2) at (5.9,-0.35) {$d_B$};
\node[enode] (e1) at (1.0,2.75) {实体};
\node[enode] (e2) at (3.35,2.75) {实体};
\node[mnode] (m1) at (7.2,2.65) {MeSH};
\node[mnode] (m2) at (8.65,1.45) {父级};
\node[mnode] (m3) at (8.65,-0.35) {兄弟};
\node[rnode] (rel) at (7.2,-1.55) {关系};

\begin{scope}[on background layer]
  \node[hedge, draw=blue!70, fill=blue!25, fit=(q)(e1)(e2), inner sep=7pt] {};
  \node[hedge, draw=cyan!70, fill=cyan!28, fit=(p1)(p2)(e1), inner sep=7pt] {};
  \node[hedge, draw=green!55!black, fill=green!25, fit=(d1)(m1)(m2), inner sep=7pt] {};
  \node[hedge, draw=orange!75!black, fill=orange!30, fit=(p2)(p3)(rel), inner sep=7pt] {};
\end{scope}

\draw[edge] (q) -- (p1);
\draw[edge] (q) -- (p2);
\draw[edge] (q) -- (p3);
\draw[edge] (p1) -- (d1);
\draw[edge] (p2) -- (d2);
\draw[edge] (p3) -- (d1);
\draw[arrow] (m1) -- (m2);
\draw[arrow] (m1) -- (m3);
\draw[edge] (rel) -- (d2);

\node[legend, anchor=west] at (9.65,1.95) {
  \textbf{超边类型}\\
  查询--概念\\
  段落--概念\\
  文档--MeSH\\
  共享候选簇\\
  层级 / 关系
};
\node[anchor=west, font=\footnotesize, text=black!62, align=left] at (9.65,-0.72) {
  扩散与中心性在每个\\问题的局部图内计算。
};
\end{tikzpicture}}
\caption{KCH-MedRank 使用的知识约束局部超图。超边连接两个以上的类型化节点，使候选段落能够从共享生物医学实体、MeSH 层级结构、文档标注和辅助关系信号中获得证据。}
\label{fig:local-hypergraph}
\end{figure}
